\documentclass{article}

\PassOptionsToPackage{numbers, sort, compress}{natbib}
\usepackage[main,final]{neurips_2025}

\usepackage[T1]{fontenc}
\usepackage[utf8]{inputenc}
\usepackage{hyperref}
\usepackage{url}
\usepackage{booktabs}
\usepackage{nicefrac}
\usepackage{microtype}
\usepackage{xcolor}
\usepackage[english,russian]{babel}


%%%

\usepackage{subcaption}
\usepackage{graphicx}
\usepackage{multirow}
\usepackage{amsmath,amssymb,amsfonts}
\usepackage{amsthm}
\usepackage{mathrsfs}
\usepackage{xcolor}
\usepackage{textcomp}
\usepackage{manyfoot}
\usepackage{algorithm}
\usepackage{algorithmicx}
\usepackage{algpseudocode}
\usepackage{listings}

\newtheorem{theorem}{Theorem} % continuous numbers
%%\newtheorem{theorem}{Theorem}[section] % sectionwise numbers
%% optional argument [theorem] produces theorem numbering sequence instead of independent numbers for Proposition
\newtheorem{proposition}[theorem]{Proposition}% 
\newtheorem{lemma}{Lemma}% 
%%\newtheorem{proposition}{Proposition} % to get separate numbers for theorem and proposition etc.

\newtheorem{example}{Example}
\newtheorem{remark}{Remark}

\newtheorem{definition}{Definition}
\newtheorem{assumption}{Assumption}

%%%


\title{MM-Algorithm for Categorical Schrödinger Bridge Matching (CSBM)}


% The \author macro works with any number of authors. There are two commands
% used to separate the names and addresses of multiple authors: \And and \AND.
%
% Using \And between authors leaves it to LaTeX to determine where to break the
% lines. Using \AND forces a line break at that point. So, if LaTeX puts 3 of 4
% authors names on the first line, and the last on the second line, try using
% \AND instead of \And before the third author name.


\author{
  Fedor Kolotilin\\
  MIPT\\
  Moscow, Russia\\
  \texttt{email@phystech.edu}\\
  \And
  Grigoriy Ksenofontov\\
  Consultant\\
  Moscow, Russia\\
  \texttt{email@phystech.edu}\\
  \AND
  Alexandr Korotin, PhD\\
  Advisor\\
  Moscow, Russia\\
  \texttt{email@phystech.edu}\\
}


\begin{document}


\maketitle

\begin{abstract}
    Задача мостов Шрёдингера играет важную роль в современном машинном обучении, позволяя решать задачи генеративного моделирования при помощи теории оптимального транспорта. Как в непрерывной так и в дискретной постановке этой задачи есть метод LightSB, сводящий задачу к оптимизации лосса (8, \cite{korotin2024light}):
    
    $$
    KL(q^* \| q_\theta) = \mathcal{L}(\theta) - \mathcal{L}^*,
    $$
    
    где
    
    $$
    \mathcal{L}(\theta) = \mathbb{E}_{p_0(x_0)} \left[ \log c_\theta(x_0) \right] - \mathbb{E}_{p_1(x_1)} \left[ \log v_\theta(x_1) \right],
    \tag{17}
    $$
    
    и $\mathcal{L}^* \in \mathbb{R}$ является константой, не зависящей от $\theta$, поэтому её можно опустить.
    
    В данной работе мы предлагаем новый метод оптимизации данной функции потерь, основанный на MM-алгоритме, который устраняет недостатки SGD: необходимость подбора гиперпараметров и шум градиентных оценок. Ключевая идея — построение суррогатной функции потерь как верхней границы исходной целевой функции. Это достигается введением вариационных параметров и применением двух границ: ELBO — для компонентов, зависящих от $v(x_1)$, и границы Бонинга \cite{boehning2010} — для компонентов, связанных с $c(x_0)$. Полученная суррогатная функция оптимизируется чередующимся алгоритмом с аналитическими обновлениями: E-шаг (вычисление оптимальных вариационных параметров при фиксированных $\theta$) и M-шаг (обновление $\theta$ в замкнутой форме при фиксированных вариационных параметрах). Итеративная процедура полностью исключает SGD, шум мини-батчей и подбор скорости обучения. Все обновления параметров становятся детерминированными и сводятся к простым аналитическим пересчетам.
\end{abstract}

\textbf{Keywords:} Categorical Schrödinger Bridge, CSBM, мост Шрёдингера, энтропийный оптимальный транспорт, дискретные диффузионные модели, MM-алгоритм, миноризация-максимизация, вариационный вывод, нижняя граница правдоподобия, ELBO, граница Бонинга, квадратичная граница, детерминированная оптимизация, аналитические обновления, непарный перевод между доменами, генеративное моделирование.

\section{Introduction}\label{sec:intro}

Categorical Schrödinger Bridge Matching (CSBM) — это метод для решения задачи построения моста Шрёдингера (Schrödinger Bridge, SB) в дискретных пространствах, предложенный Ксенофонтовым и Коротиным (2025) \cite{ksenofontov2025}. Задача SB является фундаментальным инструментом для генеративного моделирования и непарного перевода между доменами. Она заключается в поиске наиболее вероятного случайного процесса, который интерполирует между двумя заданными распределениями вероятностей.

Задачи энтропийного оптимального транспорта (EOT) и моста Шрёдингера (SB) привлекают внимание в машинном обучении благодаря приложениям в генеративном моделировании, при этом разработано множество методов для непрерывных пространств \cite{daniels2021, gushchin2023a, mokrov2024, vargas2021, chen2021, shi2023, debortoli2024, korotin2024, gushchin2024a}. Однако многие реальные данные дискретны (текст — \cite{austin2021, gat2024}; молекулярные графы — \cite{vignac2022, qin2024, luo2024}; белки — \cite{campbell2024}; векторно-квантованные представления — \cite{vandenoord2017, esser2021}). Несмотря на прогресс в дискретных диффузионных моделях \cite{hoogeboom2021, austin2021, campbell2022, lou2023, sahoo2024, campbell2024, gat2024}, существующие подходы к EOT/SB в дискретных пространствах ограничены первыми работами \cite{kim2024, ksenofontov2025}, и практические, широко применимые решатели по-прежнему отсутствуют.

\textbf{Contributions.} Our contributions can be summarized as follows:
\begin{itemize}
    \item We present a novel optimization method for the LightSB loss function based on the MM-algorithm, eliminating the drawbacks of SGD such as hyperparameter tuning and gradient noise.
    \item We demonstrate the validity of our theoretical results through empirical studies on discrete data domains.
    \item We highlight the implications of our findings for advancing the practical applicability of Schrödinger Bridge solvers in categorical settings.
\end{itemize}

\textbf{Outline.} The rest of the paper is organized as follows...

\section{Related Work}\label{sec:rw}

\textbf{Topic \#1.}
TODO

\textbf{Topic \#2.}
TODO

\section{Preliminaries}\label{sec:prelim}

\subsection{General notation}

In this section, we introduce the general notation used in the rest of the paper and the basic assumptions. 

\subsection{Assumptions} 

TODO

\section{Method}\label{sec:method}

\section{Experiments}\label{sec:exp}

To verify the theoretical estimates obtained, we conducted a detailed empirical study...

\section{Discussion}\label{sec:disc}

TODO

\section{Conclusion}\label{sec:concl}

TODO


%%%%%%%%%%%%%%%%%%%%%%%%%%%%%%%%%%%%%%%%%%%%%%%%%%%%%%%%%%%%

\bibliographystyle{unsrtnat}
\bibliography{references}

%%%%%%%%%%%%%%%%%%%%%%%%%%%%%%%%%%%%%%%%%%%%%%%%%%%%%%%%%%%%

\newpage
\appendix
\section{Appendix / supplemental material}\label{app}

\subsection{Additional experiments / Proofs of Theorems}\label{app:exp}

TODO

\end{document}